% ch13 -- Everything is Information --- But Which?
% Source documents: A180, 243, 247

\epigraph{If one sets all units to 1, ratios remain.
Pure numbers. Structures. Information.}{Matrix Book, Ch.\,16}

\section{The Fourth Reading}

The T0 theory knows three equally valid readings: Everything is energy,
everything is mass, everything is time. There is a fourth:

\textbf{Everything is information.}

What remains invariant when one sets all units to 1
($c = \hbar = \alpha = G = 1$)? \textbf{Ratios} remain.
Pure numbers. Structures. The mass ratio $m_e/m_\mu \approx 1/206{,}8$
is not energy --- it is a relation, a piece of pure information
about the geometry of the torus. \status{B}

\section{Information Without a Carrier?}

Here lies the decisive tension with A271--A273:

In the Landauer analysis, information is \emph{always} bound to a carrier.
A bit without a carrier is an abstract category, not a
physical quantity --- and has no energy. \status{B}

In the cosmological reading, the ratios $\xi, m_e/m_\mu, G/\hbar c$
are \emph{self-supporting} information --- they are not encoded in an external
medium. The universe is its own carrier.

That is not a contradiction, but a limiting-case situation: When the medium
\emph{itself} consists of information, there is no longer an independent carrier.

\begin{laierspur}
A book contains information on paper --- the carrier is the paper.
But what if the paper itself consists of letters?
Then the carrier and the content are the same. That is the situation in the
universe, if one understands it as a geometric structure: The geometry
is the carrier medium and the information simultaneously.
\end{laierspur}

\section[\texorpdfstring{$\xi$}{xi} from the Galois Algebra]{$\xi$ from the Galois Algebra}

In the FFGFT, $\xi = 4/30000$ is not an open question. It follows necessarily
from the Galois group orders of the field $\mathrm{GF}(3^n)$ and the
winding numbers of the $T^4$ torus (Doc.\,338):
\[
  \xi
  = \frac{(r_\mu/r_e)^2}{|\mathrm{GF}(9)^*|^2 \cdot 5^2 \cdot |\mathrm{GF}(27)|}
  = \frac{(12/5)^2}{8^2 \cdot 25 \cdot 27}
  = \frac{4}{30000}
\]
The three Galois factors have concrete algebraic significance: \status{K Doc.\,338}

$8 = |\mathrm{GF}(9)^*| = 3^2 - 1$ (multiplicative group of the 2nd generation),
$27 = |\mathrm{GF}(27)| = 3^3$ (field of the 3rd generation),
$5 = $ smallest prime divisor of $|\mathrm{GF}(81)^*| = 80$.

Additionally, $\xi$ is empirically anchored (Matrix Book, Ch.\,2):
$\xi = \alpha / (E_0/\mathrm{MeV})^2$.

The answer to the "why" of $\xi$ is thus algebraic, not
speculative. In Matzke's framework, the analogous question --- Why
$e_i^2 = +1$? --- remains genuinely open.

\section{Where Hyperbit and T0 Converge}

Matzke's hyperbit: \enquote{There are distinguishable things} $\to$ Geometry.

FFGFT: \enquote{There are windings on $T^4$} $\to$ Geometry, physics, masses.

Both start with pure distinguishability --- pure information ---
and end at physics. That is the deepest common point of both approaches.
\status{B}